% answers_short.tex —— 2005 年·春季北京卷·文科 · 简短答案速查
% 由 tools/make_short_answers.py 自 answers.tex 的 \gksolhead/\gkanswer 提取，
% 并与 2005/2005_数学真题_春季北京文_答案要点.md 逐题比对后生成，勿手改（改源 MD 后重跑）。

\begin{gknotice}
\textbf{使用说明}：本节只给最终答案便于快速核对，完整推导见解析版。答案为编写者独立推导，\textbf{非官方标准答案}。
\end{gknotice}

\gkpart{选择题}
\begin{center}\begin{tabular}{ccccccccc}
\toprule
题号 & 1 & 2 & 3 & 4 & 5 & 6 & 7 & 8 \\
\midrule
答案 & D & A & C & A & B & C & B & A \\
\bottomrule
\end{tabular}\end{center}

\gkpart{填空题}
\begin{itemize}[leftmargin=2.6em]
  \item[\textbf{9.}] $\frac12$
  \item[\textbf{10.}] 离心率 $\frac45$；准线 $x=\pm\frac{25}{4}$
  \item[\textbf{11.}] $\sin\theta=\frac13$，$\cos2\theta=\frac79$
  \item[\textbf{12.}] $(4+2\sqrt2)a^2$
  \item[\textbf{13.}] 一次函数 $6$ 个；二次函数 $18$ 个
  \item[\textbf{14.}] $-4<a<0$
\end{itemize}

\gkpart{解答题}
\begin{itemize}[leftmargin=2.6em]
  \item[\textbf{15.}] (1) $M=(\frac32,+\infty)$，$N=(-\infty,1]\cup[3,+\infty)$；(2) $M\cap N=[3,+\infty)$，$M\cup N=(-\infty,1]\cup(\frac32,+\infty)$
  \item[\textbf{16.}] (1) $AM/SM=3/2$；(2) $60^\circ$；(3) $V=9\sqrt3/8$
  \item[\textbf{17.}] (1) $a_n=2\cdot3^{n-1}$，$S_n=3^n-1$；(2) $b_n=3n-1$；(3) $U_{10}=425$
  \item[\textbf{18.}] (1) $\frac{x}{2}-\frac{y}{2k}=1$（$k\ne0$）；(2) $x_1x_2=4$，$y_1y_2=-4$；(3) $OM\perp ON$
  \item[\textbf{19.}] (1) 在 $v=40$ 千米/小时取得，最大车流量约为 $11.1$ 千辆/小时；(2) $25<v<64$ 千米/小时
  \item[\textbf{20.}] (1) $f(0)=0$，$f(1)=1$；(2) $k_1<k_2<k_3<k_4<k_5$；(3) $f(x_n)<x_n$（$n=1,2,3,4$）
\end{itemize}
